% SHARD-ID: RC-06H-ZETA-CONDITIONS
% SHARD-TITLE: The Zeta Conditions Imposed Factor by Factor on Graded Ring Norms: A Ledger and a Catalogue of Concrete Tensors
% SHARD-SUMMARY: Lines up ten conditions a zeta function should obey (positivity and integrality, genuine gas, rationality, functional equation, RH, unique fixed point, Galois grading with L-functions, sign structure, physical realisability, continuum) as exact algebraic statements on a graded lattice tensor and on a cMPS pair, and records which are automatic, which are constraints, and which are independent.
% SHARD-SUMMARY: Catalogue of explicit tensors with their matrices: the pair shift minus its diagonal (one zero, one pole), four-letter trace-preserving elliptic tensors (affine, projective, supersingular) with the functional equation and RH by visible mechanisms and a unique maximally mixed fixed point, a Pauli-letter family proving FE, RH and mixing mutually independent, tiny Dirichlet L-functions over F_2[x] and F_3[x] from the Horner transfer, Gauss and Kloosterman factors, a graph-cover Artin decomposition, and two cMPS (half entropy, and amplitude damping with an additive functional equation); a correction ledger and what a finite bond cannot reach.
% SHARD-KEYWORDS: zeta conditions, functional equation, Ramanujan, unique fixed point, mixing, Euler product, prime counts, elliptic curve, supersingular, Dirichlet L-function, Horner, Gauss sum, Kloosterman sum, Artin L-function, graph cover, cMPS, amplitude damping, independence, no-go
% SHARD-NOTES: none
% SHARD-SCRIPTS: scripts/zeta_conditions.py

\section{The Zeta Conditions Factor by Factor: Ledger and Catalogue}
\label{sec:zeta-conditions}

\emph{Question (TJO, 2026-09-14, night).} Line up all the conditions we want a zeta to obey
(functional equation, RH/Ramanujan, unique fixed point, \dots) and impose them factor by factor on
the cMPS ring norm, including tiny $L$-functions; the deliverable is completely concrete MPS/cMPS
that any practitioner would understand. One codex \texttt{gpt-6-astra} lane
(\texttt{notes/zeta-conditions/astra-brief.md}, report \texttt{astra-constructions.md}, 1882 lines,
13 scripts under \texttt{notes/zeta-conditions/finite/}, 461 checks re-run by the orchestrator);
the orchestrator re-derived the central constructions in \texttt{scripts/zeta\_conditions.py}
($26$ checks). Unreviewed by a second family; all rows \texttt{sketched}/\texttt{numerical}.
Conventions as in \cref{def:graded-lattice-tensor} and \cref{sec:graded-permutation}: letters
$\Kr{s}$ on $\mathbb C^{D_+|D_-}$, $\Edb$, $\Gdb$, Ramond norm $N_n = \Tr[\Gdb\Edb^n]$,
$Z(u) = \exp\sum N_nu^n/n = \det(1-u\Edb|_-)/\det(1-u\Edb|_+)$; cMPS $(Q,R)$, $\Tcm$,
$N(L) = \Tr[\Gdb e^{L\Tcm}]$, ring zeta $\det(z-\Tcm_-)/\det(z-\Tcm_+)$.

\subsection{The ledger}

\begin{observation}[Ten conditions, their exact forms, and their status]
\label{obs:zeta-condition-ledger}
\claimstatus{obs:zeta-condition-ledger}{sketched}
(C1) $N_n = \sum_w|\Tr(\Pi\Kr{w})|^2\ge0$ is automatic; integrality is a separate arithmetic constraint;
for cMPS $N(L)\ge0$ is automatic and integrality is void. (C2) Genuine gas: $a_d = \frac1d\sum_{e\mid d}
\operatorname{M\ddot ob}(d/e)N_e\in\mathbb Z_{\ge0}$, not automatic (the single letter $\operatorname{diag}(1,-1)$
has $N = 4,0,4,0,\dots$ and $a_2 = -2$); no intrinsic $a_d$ exists for continuous $L$, and no
finite constant cMPS has an atomic prime comb. (C3) Rationality is automatic; poles are net even,
zeros net odd multiplicities; in cMPS a zero eigenvalue is not invisible (it contributes a constant
to $N(L)$), and the degree difference is $(D_+-D_-)^2 = N(0)$. (C4) FE: parity-preserving
invariance of the reduced spectrum under $\lambda\mapsto q/\lambda$, giving $Z(1/(qu)) =
\pm q^{\chi/2}u^{\chi}Z(u)$; a sufficient visible mechanism is $J\Gdb = \Gdb J$, $J\Edb J^{-1} =
q\Edb^{-1}$ on the support; for cMPS the symmetry is additive, $\lambda\mapsto\kappa-\lambda$.
Correction to the reading of \cref{sec:weil-positivity}: inverse-closed letters do \emph{not}
give the FE for the raw doubled transfer ($B = \operatorname{diag}(1,2)$, $B^{-1}$: even $\{2,17/4\}$,
odd $\{5/2,5/2\}$); the reciprocal quadratics arise in the non-backtracking construction after the
Bass factors are removed. (C5) RH: $|\lambda|^2 = q$ on the retained odd spectrum; a sufficient
mechanism is a positive metric $M$ with $\Edb|_-^\dagger M\,\Edb|_- = qM$, necessary if semisimple;
Weil positivity of the rescaled trace sequence gives $|\lambda|\le\sqrt q$ and reciprocity the
equality; for cMPS $\re\lambda = \kappa/2$. (C6) Three separate requirements: a simple uncancelled
Perron root $q$; a strictly positive even $h$ with $\mathcal E^\dagger(h) = qh$ (a trace-preserving
gauge on the whole bond); a one-dimensional fixed space spanned by a density with every other
eigenvalue inside the unit circle (mixing; uniqueness alone allows periodic modes). The product
tensor $\operatorname{diag}(1,\pi)$ of \cref{sec:ring-norm-examples} has a rank-one Perron eigenmatrix
and no whole-bond gauge; the four-letter tensors below fix this without changing $Z$. (C7) A unitary
$G$-action commuting with $\Pi$ and permuting letters splits $\Edb = \oplus_\chi1\otimes E_\chi$,
$N_n = \sum_\chi\dim\chi\,N_n^\chi$, $Z = \prod L(u,\chi)^{\dim\chi}$; calling a factor an Artin
$L$-function needs an actual cover with Frobenius weights on primitive orbits; ``even = trivial,
odd = nontrivial'' is geometric input (curves), not a theorem about group actions (the Ihara cover
below has every mode even). (C8) automatic once the parity closure is fixed: the minus sign lives in
the $(+,-)$ coherences. (C9) Choi positivity per parity block plus parity covariance gives
homogeneous Kraus letters; a literal count state needs $\Tr(\Pi\Kr{w})\in\{0,1\}$; the curve
tensors are weighted norms, not point enumerations; cMPS kinetic regularity ($R_aR_b =
\pm R_bR_a$, $R_f^2 = 0$) is a further condition. (C10) Exact embedding of a lattice channel in a
cMPS needs $\Edb = e^{h\Tcm}$ with Choi conditional positivity; singular $\Edb$ (the affine tensor)
has none; Poissonisation $Q = -c/2$, $R_s = \Kr{s}$ always exists but changes FE/RH.
\end{observation}

\subsection{The catalogue}

\begin{proposition}[E0: remove a full subshift]
\label{prop:e0-pair-shift}
\claimstatus{prop:e0-pair-shift}{sketched}
Letters $\Kr{(a,b)} = \operatorname{diag}(1,[a = b])$ on $\mathbb C^{1|1}$, $a,b\in\{1,\dots,m\}$. The Ramond
ring is the indicator of the words with at least one off-diagonal pair, $N_n = m^{2n}-m^n$,
$Z = (1-mu)/(1-m^2u)$: one pole at $1/m^2$, one zero at $1/m = q^{-1/2}$ (RH by half entropy), a genuine
gas (primitive diagonal necklaces are a subset of primitive pair necklaces; $a_d = 2,5,18,57,198,661$
for $m = 2$), no FE, no whole-bond channel gauge. For $m = 2$ it equals
$Z(\mathbb A^1/\mathbb F_4)/Z(\mathbb A^1/\mathbb F_2)$ in an identified degree variable; the swap is a symmetry
with fixed letters, not an Artin cover.
\end{proposition}

\begin{proof}
$\Tr(\Pi\Kr{w}) = 1-\prod_i[a_i = b_i]$; the words are orthonormal. Orbit inclusion gives $a_d\ge0$.
\end{proof}

\begin{proposition}[Four-letter trace-preserving elliptic tensors on $\mathbb C^{1|1}$]
\label{prop:four-letter-elliptic}
\claimstatus{prop:four-letter-elliptic}{sketched-conditional}
Let $E/\mathbb F_q$ have trace $a$ and Frobenius $\pi$, $\pi\bar\pi = q$; put $r\in\{0,1\}$ (affine,
projective), $t = (q+r)/2$, $h = (q-r)/2$, and
$\Kr{0} = \operatorname{diag}(\sqrt t,\pi/\sqrt t)$, $\Kr{1} = \operatorname{diag}(0,\sqrt{t-q/t})$,
$\Kr{2} = \sqrt h\,|0\rangle\langle1|$, $\Kr{3} = \sqrt h\,|1\rangle\langle0|$. Then
$\sum\Kr{s}^\dagger\Kr{s} = q\,1$ (a trace-preserving channel after $/\sqrt q$ on the whole bond),
the even spectrum is $\{q,r\}$, the odd spectrum $\{\pi,\bar\pi\}$, and the Ramond norm is
$r+q^n-\pi^n-\bar\pi^n$: the point count of the affine ($r = 0$) or projective ($r = 1$) curve.
For $r = 1$: FE by the duality $J = \operatorname{diag}(1,-1)_{\text{even}}\oplus\text{swap}_{\text{odd}}$
with $J\Edb J^{-1} = q\Edb^{-1}$; RH because $\Edb|_-^\dagger\Edb|_- = q$; the normalised channel is
mixing with the maximally mixed state as its unique fixed point; the $H^0$ eigenvalue $1$ is a
decaying population-difference mode, not a decoupled block. Example $y^2 = x^3+x+1/\mathbb F_5$:
$a = -3$, $\pi = (-3+i\sqrt{11})/2$, $N_n = 9,27,108,675,3069,15552$. Example $y^2+y = x^3/\mathbb F_4$:
$a = -4$, $\pi = -2$, $P = (1+2u)^2$, a real double zero at $u = -1/2$, FE and RH hold, the
Artin--Schreier involution is $U = \Pi$.
\end{proposition}

\begin{proof}
Direct: the even block of $\Edb$ is $t\,1+h\,X$ with eigenvalues $t\pm h = q, r$; the odd block is
$\operatorname{diag}(\bar\pi,\pi)$; $3\cdot1+2X\mapsto3\cdot1-2X$ under $Z$-conjugation gives the duality.
Verified in \texttt{scripts/zeta\_conditions.py} (explicit words for $n\le3$) and lane script
\texttt{curve\_projective.py} ($46$ checks).
\end{proof}

\begin{proposition}[FE, RH and mixing are mutually independent]
\label{prop:fe-rh-mixing-independent}
\claimstatus{prop:fe-rh-mixing-independent}{sketched}
With Pauli letters $\sqrt{c_0}\,1,\sqrt{c_1}Z,\sqrt{c_2}X,\sqrt{c_3}Y$ on $\mathbb C^{1|1}$,
$c = \tfrac14(16+r+x+y,\ 16+r-x-y,\ 16-r+x-y,\ 16-r-x+y)$, the even spectrum is $\{16,r\}$ and the
odd $\{x,y\}$, the family is trace preserving after $/4$, and every $a_d$ is a nonnegative integer.
$(r;x,y) = (1;4,4), (1;8,2), (0;4,4), (1;3,3)$ realise (FE, RH, mixing) $=$ (Y,Y,Y), (Y,N,Y), (N,Y,Y),
(N,N,Y); tensoring with an idle even qubit keeps FE and RH and destroys uniqueness. The one
excluded combination is (FE $=$ N, RH $=$ Y) with the exact projective pair $\{1,q\}$, since RH
already makes the odd multiset reciprocal. Positivity of norms plus FE does not force RH.
\end{proposition}

\begin{proof}
The channel action of Pauli letters on $1, X, Y, Z$ is diagonal with the stated eigenvalues; orbit
embeddings give $a_d\ge0$. \texttt{independence.py} ($131$ checks), orchestrator script.
\end{proof}

\begin{proposition}[Tiny $L$-functions as MPS: Horner Dirichlet, Gauss, Kloosterman, a graph cover]
\label{prop:tiny-l-functions}
\claimstatus{prop:tiny-l-functions}{sketched}
(i) Horner MPS: bond $\mathbb F_q[x]/M$, letters $H_c$, $f\mapsto xf+c$ ($c\in\mathbb F_q$), initial vector the
residue $1$, closing row $(\chi(f))_f$; the open-chain amplitude $b_n = \chi^{\mathsf T}(\sum_cH_c)^n|1\rangle$
gives $L(u,\chi) = \sum b_nu^n$. Over $\mathbb F_2$, $M = x^3+x+1$, $\chi(x) = w = e^{2\pi i/7}$:
$L = 1+(w+w^3)u-(1+w+w^3)u^2 = (1-u)(1-\lambda u)$, $|\lambda|^2 = 2$. Over $\mathbb F_3$, $M = x^2+1$,
$\chi(1+x) = e^{2\pi i/8}$: $L = 1-\lambda u$, $\lambda = -\sqrt2+i$, $|\lambda|^2 = 3$, for the odd characters.
These are amplitudes, not norms (C1, C2 fail for an individual factor); the physical open norm counts
monics coprime to $M$; the ring presentation is the Euler product over prime polynomials with
$\chi(P)^k$ on repetitions, not the closed Horner automaton (closing $H_c$ into a trace imposes a
return condition and is wrong); the completed factor obeys the dual-character FE
$L^*(1/(qu),\chi) = -(\lambda/qu)L^*(u,\bar\chi)$. (ii) Gauss: for the quadratic character over $\mathbb F_3$
the standard convention gives $1+Gu$, $G = i\sqrt3$; a single Gauss factor is not a norm.
(iii) Kloosterman over $\mathbb F_4$: $S_n = \sum_{x\ne0}(-1)^{\operatorname{tr}(x+x^{-1})}$, $S_1 = 3$,
$L = 1+3u+4u^2 = (1-\pi u)(1-\bar\pi u)$, $\pi = (-3+i\sqrt7)/2$; the norm $4^n+S_n$ on $\mathbb C^{1|1}$
is the affine count of $y^2+xy = x^3+1$, with the sign fixed by the raw-sum convention.
(iv) Graph cover: the rose with two loops and $\mathbb Z/2$ voltages $0,1$; eight row letters on bond
dimension $8$ form a count MPS; $L_1 = 1/((1-u^2)(1-u)(1-3u))$, $L_{\mathrm{sgn}} = 1/((1-u^2)(1+3u^2))$,
every root a pole, the reduced Bass quadratics $1-4u+3u^2$ and $1+3u^2$ reciprocal at $q = 3$ (the
actual inverse-pairing mechanism); the ratio of reduced determinants $(1+3u^2)/((1-u)(1-3u))$ has a
graded $\mathbb C^{1|1}$ norm realisation. The correct tiny $\mathbb Z/2$ Artin example is a directed
voltage cover with $L(u,1) = 1/(1-3u)$, $L(u,\mathrm{sgn}) = 1/(1+u)$, not the drafted square root for the
pair shift.
\end{proposition}

\begin{proof}
Lane scripts \texttt{dirichlet\_f2.py}, \texttt{dirichlet\_f3.py}, \texttt{gauss.py}, \texttt{kloosterman.py},
\texttt{graph\_artin.py}, \texttt{z2\_shift\_cover.py}; the $\mathbb F_2$ Horner MPS re-derived by the orchestrator.
\end{proof}

\begin{proposition}[Two cMPS on $\mathbb C^{1|1}$]
\label{prop:cmps-half-and-fe}
\claimstatus{prop:cmps-half-and-fe}{sketched}
(i) Half entropy: $Q = \tfrac12\,1$, $R = \operatorname{diag}(1,0)$ give $\Tcm = \operatorname{diag}(2,1,1,1)$,
$N(L) = e^{2L}-e^{L}$, ring zeta $(z-1)/(z-2)$; a normalised finite-norm alternative
($R_f = |0\rangle\langle1|$, $R_b = \operatorname{diag}(\tfrac12,-\tfrac12)$, $Q = \operatorname{diag}(-\tfrac18,-\tfrac58)$) has
$N(L) = 1-e^{-L}$ but fails kinetic regularity; on homogeneous regular $\mathbb C^{1|1}$ data,
regularity, Lindblad normalisation, uniqueness and the exact half difference cannot all hold.
(ii) Functional equation: one lowering fermion $R = \sqrt2\,|0\rangle\langle1|$, $H = \operatorname{diag}(0,1)$,
$Q = -iH-\tfrac12R^\dagger R$: normalised spectrum even $\{0,-2\}$, odd $\{-1\pm i\}$, unique pure
vacuum fixed point; in the growth convention $N(L) = |1-e^{(1+i)L}|^2 = 1+e^{2L}-2e^L\cos L$ and
ring zeta $((z-1)^2+1)/(z(z-2))$ with the additive FE $\mathcal Z(2-z) = \mathcal Z(z)$ and both zeros on
$\re z = 1$: FE, RH and the unique fixed point by one visible mechanism (odd damping is half the
population decay, the frequency comes from $H$). The ring vector has only its vacuum amplitude:
the zeta data sit in the length-dependent vacuum normalisation, and no nontrivial entanglement
spectrum is produced.
\end{proposition}

\begin{proof}
Direct diagonalisation of $\Tcm$; \texttt{cmps\_half.py}, \texttt{cmps\_fe.py}, orchestrator script.
\end{proof}

\begin{observation}[What factor by factor reaches, and what a finite bond cannot]
\label{obs:factor-by-factor-limits}
\claimstatus{obs:factor-by-factor-limits}{sketched-conditional}
A practitioner can build tonight: E0 (a zero from destructive interference with a subshift), the
four-letter elliptic channels (RH, then FE on completion, with a unique mixing fixed point), the
Pauli family (break RH keeping FE and mixing; add an idle qubit to break uniqueness without moving
any divisor point), the Horner tensors (tiny Dirichlet polynomials), the graph cover (an Artin
decomposition whose nontrivial factors have poles), the damping cMPS (additive FE, uniform odd
rate, unique vacuum). A finite bond can have infinitely many Euler primes (E0 does) but cannot
produce a nonrational $Z$, an atomic prime comb, or infinitely many distinct divisor points:
$u = q^{-s}$ gives periodic copies of finitely many points, not Riemann's zeros. The no-go of
\cref{sec:ring-norm-bosonic-nogo} concerns plain power traces; graded norms evade it by negative net
coherence multiplicities. A single simple pole at $1/q$ is incompatible with the uncompleted
self-reciprocal FE: the completion must be declared. Factor by factor constructs and tests local
spectra, characters, primes, normalisations and dualities; it does not make an infinite product
converge, supply the continuation or the archimedean factor, select an entanglement spectrum, or
prove RH for an unidentified infinite object. The corrections to the orchestrator's drafts are in
the lane's ledger (D1--D8): inverse pairing alone gives no FE; the swap $L$-function was not an
Artin factor; unique stationarity, gauge and mixing are three conditions; the affine rung needs its
own tensor.
\end{observation}

\begin{numerical}[Zeta conditions: one lane, two suites]
\label{num:zeta-conditions}
\claimstatus{num:zeta-conditions}{numerical}
Lane: $461$ checks over $13$ scripts (\texttt{notes/zeta-conditions/finite/run\_all.py}, re-run by the
orchestrator, all pass; every matrix printed in the report's generated catalogue). Orchestrator:
\texttt{scripts/zeta\_conditions.py} (\texttt{outputs/zeta\_conditions.txt}), $26$ checks: E0 counts, gas
and zeta; the $\mathbb F_5$ curve count, the four-letter TP tensor, its spectra, Ramond norm by explicit
words, the duality $J$, the unitary odd block, the unique mixing fixed point; the four Pauli
families with their FE/RH/gas verdicts; the FE cMPS spectrum, $N(L)$ at six lengths, and the
additive FE; the $\mathbb F_2$ Horner $L$-function and $|\lambda|^2 = 2$; the $\mathbb F_4$ supersingular count.
\end{numerical}
